% ============================================================
% metodologia.tex — Metodología CRISP-DM aplicada al caso
% ============================================================
\section{Metodología}

Se adoptó la metodología CRISP-DM (\textit{Cross-Industry Standard Process
for Data Mining}) \cite{wirth2000crisp}, que estructura el desarrollo en
seis fases iterativas.

\subsection{Comprensión del Negocio}
El problema se formula como clasificación binaria supervisada: dada una
transacción con características $\mathbf{x} \in \mathbb{R}^{30}$, predecir
$y \in \{0, 1\}$ donde $y=1$ indica fraude. La métrica de negocio
prioritaria es el \textit{Recall} de la clase fraude, pues un falso
negativo (fraude no detectado) implica pérdida económica directa.

\subsection{Comprensión y Preparación de los Datos}
El dataset contiene 284,807 transacciones con 30 características:
28 componentes PCA ($V_1$–$V_{28}$), \texttt{Time} y \texttt{Amount}.
La variable objetivo \texttt{Class} presenta una distribución extremadamente
desbalanceada: 99.83\% transacciones legítimas vs 0.17\% fraudulentas.

Las columnas \texttt{Time} y \texttt{Amount} se normalizaron con
\texttt{StandardScaler}. El conjunto se dividió 80/20 con estratificación
para preservar la proporción de clases.

\subsection{Manejo del Desbalance — SMOTE}
Se aplicó \textit{Synthetic Minority Over-sampling Technique} (SMOTE)
\cite{chawla2002smote} exclusivamente sobre el conjunto de entrenamiento
para evitar data leakage. SMOTE genera instancias sintéticas de la clase
minoritaria interpolando entre vecinos cercanos en el espacio de
características.

\subsection{Modelos Evaluados}

\subsubsection{Regresión Logística}
Modelo lineal base. Dado el vector de características $\mathbf{x}$:
\begin{equation}
P(y=1|\mathbf{x}) = \sigma(\mathbf{w}^T\mathbf{x} + b) = \frac{1}{1+e^{-(\mathbf{w}^T\mathbf{x}+b)}}
\end{equation}

\subsubsection{Random Forest}
Ensemble de $T$ árboles de decisión. La predicción final agrega por
mayoría de votos:
\begin{equation}
\hat{y} = \text{mode}\{h_t(\mathbf{x})\}_{t=1}^{T}
\end{equation}

\subsubsection{XGBoost}
Gradient boosting optimizado. Minimiza una función objetivo regularizada:
\begin{equation}
\mathcal{L} = \sum_i l(y_i, \hat{y}_i) + \sum_k \Omega(f_k)
\end{equation}

\subsection{Métricas de Evaluación}
Dado el desbalance, se descarta la exactitud como métrica primaria.
Se utilizan:
\begin{itemize}
  \item \textbf{ROC-AUC}: área bajo la curva ROC
  \item \textbf{Precision}: $TP / (TP + FP)$
  \item \textbf{Recall}: $TP / (TP + FN)$ — métrica principal
  \item \textbf{F1-Score}: media armónica de Precision y Recall
\end{itemize}
